\begin{figure}[htbp]
\centering
\resizebox{\textwidth}{!}{%
\begin{tikzpicture}[font=\sffamily]

% =========================
% Color Definitions
% =========================
\colorlet{col1}{blue!18}
\colorlet{col2}{green!18}
\colorlet{col3}{teal!22}
\colorlet{col4}{gray!20}

\colorlet{line1}{blue!60!black}
\colorlet{line2}{green!50!black}
\colorlet{line3}{teal!70!black}
\colorlet{line4}{gray!70!black}

% =========================
% Pyramid Geometry
% =========================
\def\W{9.5}   % Slightly widened base to easily fit the \normalsize text
\def\H{9.6}   % Total height
\def\CX{-3.0} % Center X of pyramid to balance the boxes on the right

\coordinate (BL) at (\CX-\W/2, 0);
\coordinate (BR) at (\CX+\W/2, 0);
\coordinate (TOP) at (\CX, \H);

% Y-coordinates of slice boundaries
\def\yA{2.4}
\def\yB{4.8}
\def\yC{7.2}

% Calculate coordinates on the edges
\coordinate (L1) at ($(BL)!\yA/\H!(TOP)$);
\coordinate (R1) at ($(BR)!\yA/\H!(TOP)$);
\coordinate (L2) at ($(BL)!\yB/\H!(TOP)$);
\coordinate (R2) at ($(BR)!\yB/\H!(TOP)$);
\coordinate (L3) at ($(BL)!\yC/\H!(TOP)$);
\coordinate (R3) at ($(BR)!\yC/\H!(TOP)$);

% =========================
% Filled layers & Borders
% =========================
\fill[col1] (BL) -- (BR) -- (R1) -- (L1) -- cycle;
\fill[col2] (L1) -- (R1) -- (R2) -- (L2) -- cycle;
\fill[col3] (L2) -- (R2) -- (R3) -- (L3) -- cycle;
\fill[col4] (L3) -- (R3) -- (TOP) -- cycle;

% Internal layer borders
\draw[white, line width=1.5pt] (L1) -- (R1);
\draw[white, line width=1.5pt] (L2) -- (R2);
\draw[white, line width=1.5pt] (L3) -- (R3);

% Outer pyramid border
\draw[black!60, line width=1pt] (BL) -- (BR) -- (TOP) -- cycle;

% =========================
% Box Styles
% =========================
\tikzset{
    infobox/.style={
        draw,
        rounded corners=2.5pt,
        line width=0.8pt,
        fill=white,
        align=left,
        text width=4.5cm,
        inner sep=6pt,
        font=\footnotesize % Reduced to standard diagram font size
    }
}

% =========================
% Edge Connectors & Boxes
% =========================
\coordinate (M1) at ($(BR)!0.5!(R1)$);
\coordinate (M2) at ($(R1)!0.5!(R2)$);
\coordinate (M3) at ($(R2)!0.5!(R3)$);
\coordinate (M4) at ($(R3)!0.5!(TOP)$);

\def\boxX{2.2} % Left edge of all boxes

% Box 1 (Mean)
\node[infobox, draw=line1] (box1) at (\boxX, 1.2) [anchor=west] {
    \textbf{Target domain:} attenuation\\[0.1cm]
    \textbf{Relative role:} useful first-order signal
};
\draw[line1, line width=0.9pt] (M1) -- (box1.west);

% Box 2 (Ring)
\node[infobox, draw=line2] (box2) at (\boxX, 3.6) [anchor=west] {
    \textbf{Target domain:} attenuation + coherence\\[0.1cm]
    \textbf{Relative role:} broadest explanatory improvement
};
\draw[line2, line width=0.9pt] (M2) -- (box2.west);

% Box 3 (Water)
\node[infobox, draw=line3] (box3) at (\boxX, 6.0) [anchor=west] {
    \textbf{Target domain:} coherence-focused\\[0.1cm]
    \textbf{Relative role:} strong but narrower contribution
};
\draw[line3, line width=0.9pt] (M3) -- (box3.west);

% Box 4 (Morphology)
\node[infobox, draw=line4] (box4) at (\boxX, 8.4) [anchor=west] {
    \textbf{Target domain:} structural support\\[0.1cm]
    \textbf{Relative role:} secondary explanatory layer
};
\draw[line4, line width=0.9pt] (M4) -- (box4.west);

% =========================
% Internal Pyramid Labels
% =========================
\node[font=\normalsize\bfseries] at (\CX, 1.2) {Mean land-cover contrast};
\node[font=\normalsize\bfseries] at (\CX, 3.6) {Ring-gradient structure};
\node[font=\normalsize\bfseries] at (\CX, 6.0) {Water context};
% Morphology is given a smaller font and lowered slightly to safely clear the borders
\node[font=\scriptsize\bfseries] at (\CX, 8.1) {Morphology};

% =========================
% Grouping Brace
% =========================
\draw[decorate,decoration={brace,amplitude=5pt,mirror},line width=1pt,black!70]
(7.0, 0.4) -- (7.0, 4.4)
node[midway, right=10pt, align=left, font=\small\bfseries, text width=2.8cm, color=black!85]
{Dominant\\[2pt]land-cover\\[2pt]explanatory\\[2pt]base};

\end{tikzpicture}%
}
\caption{Conceptual synthesis of the layered explanatory hierarchy identified in the Bangkok park--outside PM$_{2.5}$ wavelet analysis. Mean land-cover contrast provided a useful first-order signal, ring-gradient structure provided the broadest explanatory improvement across attenuation and coherence, water context contributed strong but narrower coherence-focused associations, and morphology acted as a secondary structural layer rather than a dominant control.}
\label{fig:hierarchy_pyramid}
\end{figure}